\documentclass[10pt,a4paper]{amsart}
\usepackage{amsmath,amssymb,amsthm,graphicx}
\usepackage{verbatim}
\newenvironment{reference}{\comment}{\endcomment}
%\newenvironment{reference}{}{}
\newenvironment{slogan}{\comment}{\endcomment}
\newenvironment{history}{\comment}{\endcomment}

% For commutative diagrams we use Xy-pic
\usepackage[all]{xy}

% We use 2cell for 2-commutative diagrams.
\xyoption{2cell}
% Theorem environments.
%
\theoremstyle{plain}
\newtheorem{theorem}[subsection]{Theorem}
\newtheorem{proposition}[subsection]{Proposition}
\newtheorem{lemma}[subsection]{Lemma}

\theoremstyle{definition}
\newtheorem{definition}[subsection]{Definition}
\newtheorem{example}[subsection]{Example}
\newtheorem{exercise}[subsection]{Exercise}
\newtheorem{situation}[subsection]{Situation}

\theoremstyle{remark}
\newtheorem{remark}[subsection]{Remark}
\newtheorem{remarks}[subsection]{Remarks}

\numberwithin{equation}{subsection}

% Macros
%
\def\lim{\mathop{\mathrm{lim}}\nolimits}
\def\colim{\mathop{\mathrm{colim}}\nolimits}
\def\Spec{\mathop{\mathrm{Spec}}}
\def\Hom{\mathop{\mathrm{Hom}}\nolimits}
\def\Ext{\mathop{\mathrm{Ext}}\nolimits}
\def\SheafHom{\mathop{\mathcal{H}\!\mathit{om}}\nolimits}
\def\SheafExt{\mathop{\mathcal{E}\!\mathit{xt}}\nolimits}
\def\Sch{\mathit{Sch}}
\def\Mor{\mathop{\mathrm{Mor}}\nolimits}
\def\Ob{\mathop{\mathrm{Ob}}\nolimits}
\def\Sh{\mathop{\mathit{Sh}}\nolimits}
\def\NL{\mathop{N\!L}\nolimits}
\def\CH{\mathop{\mathrm{CH}}\nolimits}
\def\proetale{{pro\text{-}\acute{e}tale}}
\def\etale{{\acute{e}tale}}
\def\QCoh{\mathit{QCoh}}
\def\Ker{\mathop{\mathrm{Ker}}}
\def\Im{\mathop{\mathrm{Im}}}
\def\Coker{\mathop{\mathrm{Coker}}}
\def\Coim{\mathop{\mathrm{Coim}}}

% Boxtimes
%
\DeclareMathSymbol{\boxtimes}{\mathbin}{AMSa}{"02}

%
% Macros for moduli stacks/spaces
%
\def\QCohstack{\mathcal{QC}\!\mathit{oh}}
\def\Cohstack{\mathcal{C}\!\mathit{oh}}
\def\Spacesstack{\mathcal{S}\!\mathit{paces}}
\def\Quotfunctor{\mathrm{Quot}}
\def\Hilbfunctor{\mathrm{Hilb}}
\def\Curvesstack{\mathcal{C}\!\mathit{urves}}
\def\Polarizedstack{\mathcal{P}\!\mathit{olarized}}
\def\Complexesstack{\mathcal{C}\!\mathit{omplexes}}
% \Pic is the operator that assigns to X its picard group, usage \Pic(X)
% \Picardstack_{X/B} denotes the Picard stack of X over B
% \Picardfunctor_{X/B} denotes the Picard functor of X over B
\def\Pic{\mathop{\mathrm{Pic}}\nolimits}
\def\Picardstack{\mathcal{P}\!\mathit{ic}}
\def\Picardfunctor{\mathrm{Pic}}
\def\Deformationcategory{\mathcal{D}\!\mathit{ef}}
\usepackage{cleveref}
\crefname{originalproposition}{original proposition}{original propositions}
\Crefname{originalproposition}{Original proposition}{Original propositions}
\makeatletter
\newlabel{proposition-finite-higher-direct-image-zero}{{03QP}{source}{}{sourcepointer-1}{}}
\newlabel{proposition-finite-higher-direct-image-zero@cref}{{[originalproposition][0][]03QP}{[1][0][]source}{}{}{}}
\makeatother
\crefname{originallemma}{original lemma}{original lemmas}
\Crefname{originallemma}{Original lemma}{Original lemmas}
\makeatletter
\newlabel{lemma-jshriek-open}{{0F70}{source}{}{sourcepointer-2}{}}
\newlabel{lemma-jshriek-open@cref}{{[originallemma][0][]0F70}{[1][0][]source}{}{}{}}
\makeatother
\crefname{originallemma}{original lemma}{original lemmas}
\Crefname{originallemma}{Original lemma}{Original lemmas}
\makeatletter
\newlabel{lemma-finite-pushforward-commutes-with-base-change}{{0959}{source}{}{sourcepointer-3}{}}
\newlabel{lemma-finite-pushforward-commutes-with-base-change@cref}{{[originallemma][0][]0959}{[1][0][]source}{}{}{}}
\makeatother
\title{Base change compatibility of extension by zero and pushforward along a finite morphism}
\author{The Stacks Project Authors}
\date{stacks-project:etale-cohomology.tex:local-snapshot-20261011; GNU FDL 1.2 or later}
\begin{document}
\maketitle
\noindent\textit{Editorial note.} All delivered mathematical text is the author's own, copied from the pinned Stacks Project chapter file etale-cohomology.tex; no statement or proof step was rewritten, reordered or removed.\par\medskip
\noindent\textit{Editorial note.} The proof is the author's complete argument for the identification of the two functors; the closing sentence omits only the verification of functoriality in the sheaf, which is a routine compatibility, and it is preserved verbatim in the counted span.\par\medskip
\noindent\textit{Editorial note.} Every labelled reference occurring inside the counted spans is carried in assembly.external\_references as a non-load-bearing navigation pointer (same-chapter ones at their definition site, cross-chapter ones at the invocation inside the counted span), so no load-bearing mathematics is replaced by a stub. No bibliography citation occurs inside the counted spans.\par\medskip
\noindent\textit{Editorial note.} The checked chapter list for same-name collisions is the set of stacks chapters staged in this repository: chow, cohomology, divisors, etale-cohomology, flat, more-algebra, more-morphisms, morphisms, spaces-more-morphisms, varieties together with this card's chapter etale-cohomology.\par\medskip
\noindent\textit{Editorial note.} The preamble is the chapter's own native preamble excerpt already used for this source: the Xy-pic setup (native preamble lines 14-18) and the native theorem declarations with the mathematical macros (native preamble lines 154-217).\par\medskip
\noindent\textit{Allowed standard prerequisites (editorial).}\begin{itemize}
\item Vanishing of the higher direct images of a finite morphism. Same chapter, Proposition proposition-finite-higher-direct-image-zero. Computes the stalks of the pushforward as a finite direct sum of stalks.
\item Description of extension by zero and its stalks. Same chapter, Lemma lemma-jshriek-open. Characterizes extension by zero by the vanishing of its stalks outside the open and by its restriction.
\item Pushforward commutes with base change for finite morphisms. Same chapter, Lemma lemma-finite-pushforward-commutes-with-base-change. Identifies the restriction of the pushforward with the pushforward of the restriction.
\end{itemize}\medskip
\ifdefined\hypertarget\hypertarget{sourcepointer-1}{}\fi
\noindent\textit{Source reference.} Original proposition 03QP: sources/stacks-etale-cohomology/pinned-original-source.tex, source lines 8245--8245. Navigation pointer invoked at proof line 11634; the counted span contains the pointer only and the cited result is not reproduced inside it, so the pointer is registered as a non-load-bearing navigation stub in the established form for Stacks material. No bibliography citation is involved.\par\medskip
\ifdefined\hypertarget\hypertarget{sourcepointer-2}{}\fi
\noindent\textit{Source reference.} Original lemma 0F70: sources/stacks-etale-cohomology/pinned-original-source.tex, source lines 11468--11468. Navigation pointer invoked at proof line 11636; the counted span contains the pointer only and the cited result is not reproduced inside it, so the pointer is registered as a non-load-bearing navigation stub in the established form for Stacks material. No bibliography citation is involved.\par\medskip
\ifdefined\hypertarget\hypertarget{sourcepointer-3}{}\fi
\noindent\textit{Source reference.} Original lemma 0959: sources/stacks-etale-cohomology/pinned-original-source.tex, source lines 8292--8292. Navigation pointer invoked at proof line 11642; the counted span contains the pointer only and the cited result is not reproduced inside it, so the pointer is registered as a non-load-bearing navigation stub in the established form for Stacks material. No bibliography citation is involved.\par\medskip
\begin{lemma}
\label{lemma-compatible-shriek-push-finite}
Consider a cartesian diagram of schemes
$$
\xymatrix{
U \ar[d]_g \ar[r]_{j'} & X \ar[d]^f \\
V \ar[r]^j & Y
}
$$
where $f$ is finite and $j$ is an open immersion.
Then $f_* \circ j'_! = j_! \circ g_*$ as functors
$\textit{Ab}(U_\etale) \to \textit{Ab}(Y_\etale)$.
\end{lemma}

\begin{proof}
Let $\mathcal{F}$ be an object of $\textit{Ab}(U_\etale)$.
Let $\overline{y}$ be a geometric point of $Y$ not contained
in the open $V$. Then
$$
(f_*j'_!\mathcal{F})_{\overline{y}} =
\bigoplus\nolimits_{\overline{x},\ f(\overline{x}) = \overline{y}}
(j'_!\mathcal{F})_{\overline{x}} = 0
$$
by Proposition \ref{proposition-finite-higher-direct-image-zero} and
because the stalk of $j'_!\mathcal{F}$ at $\overline{x} \not \in U$
are zero by Lemma \ref{lemma-jshriek-open}. On the other hand, we have
$$
j^{-1}f_*j'_!\mathcal{F} =
g_*(j')^{-1}j'_!\mathcal{F} =
g_*\mathcal{F}
$$
by Lemmas \ref{lemma-finite-pushforward-commutes-with-base-change} and
\ref{lemma-jshriek-open}.
Hence by the characterization of
$j_!$ in Lemma \ref{lemma-jshriek-open} we see that
$f_*j'_!\mathcal{F} = j_!g_*\mathcal{F}$.
We omit the verification that this identification
is functorial in $\mathcal{F}$.
\end{proof}

\end{document}
